\section{Twins, vertices outside the packing, and a shrink rule}\label{app:structure}\label{sec:structure}
By Lemma~\ref{lem:decomp} the whole difficulty of $\ell>0$ is the set $U$ of at most $\ell$ hard edits. This appendix proves what can be said about the vertices outside $V(H)$ for pair-maximal packings $H$. The results are not used in the main text.

\paragraph{Conventions.} Partitions, clusters, costs and optima are as in Section~\ref{sec:prelim}; $\Pi(v)$ is the cluster of $v$. Unless said otherwise, $H$ is pair-maximal and $R=V\setminus V(H)$; every pair with an end in $R$ is hard.

\subsection{Twins}
\begin{lemma}[Twins, strict form]\label{lem:twin}
Let $u\ne v$ be adjacent with $N[u]=N[v]$. For every partition $\Pi$ with $\Pi(u)\ne\Pi(v)$, one of the two partitions obtained by moving $u$ into $\Pi(v)$, respectively $v$ into $\Pi(u)$, has cost smaller than $\Pi$ by at least $1$; more precisely the two cost changes sum to exactly $-2$. Hence $u,v$ lie in one cluster in \emph{every} optimum, with no hypothesis on $H$.
\end{lemma}
\begin{proof}
For a vertex $y$ and a set $Y\not\ni y$ let $f(y,Y)=|Y\setminus N(y)|-|Y\cap N(y)|$. Put $C_u=\Pi(u)\setminus\{u\}$ and $C_v=\Pi(v)\setminus\{v\}$ (possibly empty); $C_u\not\ni v$, $C_v\not\ni u$, and neither contains $u$ or $v$. Moving $u$ out of $\Pi(u)$ changes the status of the pairs $uz$, $z\in C_u$: before, $uz$ is an edit iff it is a non-edge; after, iff it is an edge; the change is $-f(u,C_u)$. Moving $u$ into $\Pi(v)=C_v\cup\{v\}$ changes, for $z\in C_v\cup\{v\}$, the status from ``edit iff edge'' to ``edit iff non-edge''; the change is $f(u,C_v\cup\{v\})=f(u,C_v)-1$ because $uv\in E$. All other pairs keep their status. So
\[\delta_1=-f(u,C_u)+f(u,C_v)-1,\qquad \delta_2=-f(v,C_v)+f(v,C_u)-1 .\]
Since $N(u)\setminus\{v\}=N(v)\setminus\{u\}$ we have $f(u,Y)=f(v,Y)$ for all $Y\subseteq V\setminus\{u,v\}$, hence $\delta_1+\delta_2=-2$. Thus $\min(\delta_1,\delta_2)\le-1$.
\end{proof}
This is the critical-clique property used in kernelizations of Cluster Editing~\cite{FLRS07,Guo09} (it is also recalled in~\cite{BFK18} when motivating Reduction Rule~6.2); the argument above is self-contained and shows that it is strict.

\subsection{Viability of a hard edit}
\begin{lemma}[Viability of a hard edit]\label{lem:viab}
Let $S$ be any solution, $U$ its hard part and $uv\in U$. Let $\mathrm{Bd}(uv)$ be the set of $z\notin Q(u)\cup Q(v)\cup\{u,v\}$ such that $\{u,v,z\}$ is an induced $P_3$ of $G\sd\{uv\}$. Then for every $z\in\mathrm{Bd}(uv)$ one of $uz,vz$ lies in $U$, and $|\mathrm{Bd}(uv)|\le|U|-1$.
\end{lemma}
\begin{proof}
All three pairs of $\{u,v,z\}$ are hard ($z\notin Q(u)\cup Q(v)$). If neither $uz$ nor $vz$ is in $S$, the triple has in $G\sd S$ the same pattern as in $G\sd\{uv\}$, an induced $P_3$. Distinct $z$ give distinct pairs, all different from $uv$.
\end{proof}
This gives a kernel-style condition (deleting $uv$ needs fewer than $\ell$ common neighbours outside $Q(u)\cup Q(v)$) but no bound on the number of candidate pairs: all pendant edges are viable.

\subsection{The vertices outside \texorpdfstring{$V(H)$}{V(H)}}
\begin{lemma}[Structure of $R$]\label{lem:R}
Let $H$ be pair-maximal. For $r,s\in R$: if $rs\in E$ then $N[r]=N[s]$; if $rs\notin E$ then $N(r)\cap N(s)=\emptyset$. Consequently:
\begin{enumerate}
\item every component $X$ of $G[R]$ (a \emph{class}) is a clique all of whose vertices have the same closed neighbourhood $N[X]$ (they are true twins);
\item $A_X:=N(X)\setminus X$ is contained in $V(H)$;
\item the sets $A_X$, over the classes $X$, are pairwise disjoint.
\end{enumerate}
\end{lemma}
\begin{proof}
Every pair with an end in $R$ is hard. Let $rs\in E$ and $z\ne r,s$ be adjacent to exactly one of $r,s$. Then $\{r,s,z\}$ has exactly two edges, i.e.\ is an induced $P_3$ with all three pairs hard, contradicting pair-maximality; so $N[r]=N[s]$. If $rs\notin E$ and $z$ is a common neighbour, then $r\,z\,s$ is an induced $P_3$ whose pairs $rz,zs,rs$ all have an end in $R$, again a contradiction.

(1) Adjacent vertices of $R$ have equal closed neighbourhoods; along a connected component all closed neighbourhoods coincide, and then any two vertices $r,t$ of the component satisfy $r\in N[r]=N[t]$, so they are adjacent. (2) A neighbour of $x\in X$ lying in $R$ is in the component of $x$, hence in $X$. (3) If $a\in A_X\cap A_Y$ for classes $X\ne Y$, then $a$ is a common neighbour of some $x\in X$ and $y\in Y$, which are non-adjacent (different components of $G[R]$); this contradicts the second clause.
\end{proof}
For $x\in X$ and $a\in V(H)$ we thus have $xa\in E$ iff $a\in A_X$; in particular vertices of one class are twins towards $V(H)$ as well. (Lemma~\ref{lem:R} fails for vertex-maximal packings, see Proposition~\ref{prop:vertexmax}, whose instance also shows that Theorem~\ref{thm:C1} fails without pair-maximality.)

\begin{lemma}[Exchange formula]\label{lem:exch}
Let $\Pi$ be any partition, $X$ a class with $X\subseteq C$ for a cluster $C$, $x=|X|$ and $a=|A_X\cap C|$. Let $\Pi'$ be obtained by moving $X$ out of $C$ into a cluster of its own. Then $\mathrm{cost}(\Pi')-\mathrm{cost}(\Pi)=x\,(x+2a-|C|)$.
\end{lemma}
\begin{proof}
Only the pairs between $X$ and $C\setminus X$ change status. By Lemma~\ref{lem:R}, a vertex of $X$ has exactly $a$ neighbours in $C\setminus X$, namely $A_X\cap C$. Before the move such a pair is an edit iff it is a non-edge: $x(|C|-x-a)$ edits. After the move it is an edit iff it is an edge: $xa$ edits. The difference is $xa-x(|C|-x-a)=x(x+2a-|C|)$.
\end{proof}

\begin{theorem}[C1: no edit inside $R$]\label{thm:C1}
Let $H$ be pair-maximal. In \emph{every} optimum $\Pi$, every cluster meets $R$ in the empty set or in exactly one whole class. Hence no edit of $\Pi$ has both endpoints in $R$, and for every class $X$ the cluster $\Pi(X)$ equals $X\cup T_X$ with $T_X:=\Pi(X)\cap V(H)$.
\end{theorem}
\begin{proof}
By Lemma~\ref{lem:R}(1) and Lemma~\ref{lem:twin}, every class lies inside one cluster (a class of size $1$ trivially). Suppose a cluster $C$ contains two distinct classes $X_1,X_2$. Put $x_i=|X_i|\ge1$ and $a_i=|A_{X_i}\cap C|$. The sets $X_1,X_2\subseteq R$ and $A_{X_1}\cap C,A_{X_2}\cap C\subseteq V(H)$ are pairwise disjoint (Lemma~\ref{lem:R}(2),(3)), so $|C|\ge x_1+x_2+a_1+a_2$. Let $\Delta_i$ be the cost change of moving $X_i$ out of $C$. Since $\Pi$ is an optimum, $\Delta_1\ge0$ and $\Delta_2\ge0$, i.e.\ by Lemma~\ref{lem:exch} $|C|\le x_i+2a_i$ for $i=1,2$. With the lower bound on $|C|$: $x_2+a_2\le a_1$ and $x_1+a_1\le a_2$. Adding gives $x_1+x_2\le0$, absurd. So no cluster contains two classes.

An edit with both ends in $R$ joins two vertices $r,s$. If they lie in one class, they are adjacent and (by the above) in one cluster, so $rs$ is no edit. If they lie in different classes they are non-adjacent, and $rs$ is an edit iff they share a cluster, which was excluded.
\end{proof}
The strictness of the inequality is what makes the statement hold for every optimum (no ties, and no iteration or termination argument).

\subsection{Consequences}\label{sec:conseq}
Let $(G,H,\ell)$ be a \emph{YES instance} with $H$ pair-maximal, $\Pi$ any optimum, $S$ its edit set, so $|S|=\OPT\le|H|+\ell$. Let $U$ be the hard edits. By Lemma~\ref{lem:decomp}, $|S|-|H|=\sum_i(|x_i|-1)+|U|$, so $|U|\le\ell$.  Write $q=\max_u|Q(u)|\le 2c$.

\begin{corollary}[Normal form]\label{cor:nf}
For every class $X$ let $T_X=\Pi(X)\cap V(H)$.
\begin{enumerate}
\item[(N1)] \emph{Types of hard edits.} Every hard edit is either an \emph{HH-edit} (both ends in $V(H)$) or an \emph{RH-edit} $\{x,a\}$, $x\in X$ a class, $a\in V(H)$; the latter is a deletion if $a\in A_X\setminus T_X$ and an insertion if $a\in T_X\setminus A_X$, and $\{x,a\}$ is edited iff $\{x',a\}$ is, for all $x,x'\in X$. With $h$ the number of HH-edits,
$|U|=h+\sum_X|X|\cdot|A_X\triangle T_X|\le\ell-\sum_i(|x_i|-1)\le \ell.$
\item[(N2)] \emph{Frozen classes.} If $|X|>\ell$, then $\Pi(X)=X\cup A_X$.
\item[(N3)] \emph{Attachment.} Every $a\in V(H)$ lies in $A_X$ for at most one class $X$. If $z,z'\in A_X$ are non-adjacent then $zz'$ is soft. If $a\in V(H)\setminus A_X$, $z\in A_X$ and $az\in E$ then $az$ is soft; hence $|N(a)\cap A_X|\le|Q(a)|\le q$.
\item[(N4)] \emph{Shrink rule.} If $|X|\ge\ell+2$ and $x\in X$, then $(G,H,\ell)$ is a YES instance iff $(G-x,H,\ell)$ is. This holds for $H$ pair-maximal without the YES hypothesis. If $A_X=\emptyset$ then $X$ is a component of $G$ and may be deleted without changing $\OPT$.
\item[(N5)] \emph{Size.} For a pair-maximal $H$ and after exhausting the two rules in (N4), every class has at most $\ell+1$ vertices and at most $|V(H)|$ classes have $A_X\ne\emptyset$, so $|R|\le(\ell+1)|V(H)|$ and $n\le(\ell+2)|V(H)|\le 3(\ell+2)|H|$. ($|V(H)|$ is not bounded in terms of $\ell$, so this is not a kernel; see Remark~\ref{rem:ext}.)
\end{enumerate}
\end{corollary}
\begin{proof}
(N1) Pairs with an end in $R$ are hard; by Theorem~\ref{thm:C1} none has both ends in $R$. The vertices of $X$ are twins towards $V(H)$ and lie in one cluster, which gives the all-or-none statement; $xa$ is an edit iff exactly one of $a\in T_X$, $a\in A_X$ holds. Counting gives the identity; the bound is the first paragraph of this subsection.

(N2) If $\Pi(X)\ne X\cup A_X$ then, since $\Pi(X)=X\cup T_X$ (Theorem~\ref{thm:C1}), $A_X\triangle T_X\neq\emptyset$ and (N1) gives $|U|\ge|X|>\ell$, contradicting $|U|\le\ell$.

(N3) The first claim is Lemma~\ref{lem:R}(3). Let $x\in X$. If $z,z'\in A_X$ are non-adjacent, $z\,x\,z'$ is an induced $P_3$ with pairs $zx,xz'$ hard, so $zz'$ must be soft (pair-maximality). If $a\in V(H)\setminus A_X$ and $z\in A_X$ with $az\in E$, then $xa\notin E$ and $x\,z\,a$ is an induced $P_3$ with $xz,xa$ hard, so $az$ is soft, i.e.\ $z\in Q(a)$.

(N4) $x\in R$, so $H$ is still a packing of $G-x$, still pair-maximal (an induced $P_3$ of $G-x$ is one of $G$), and $X\setminus\{x\}$ is a class of $G-x$ with the same $A_X$. ($\Rightarrow$) Restricting a partition of $G$ to $V\setminus x$ does not increase its cost, so $\OPT(G-x)\le\OPT(G)\le|H|+\ell$. ($\Leftarrow$) Let $\Pi$ be an optimum of $G-x$ with cost $\le|H|+\ell$. Now $|X\setminus x|\ge\ell+1>\ell$, so (N2) in $G-x$ gives $\Pi(X\setminus x)=(X\setminus x)\cup A_X$. Add $x$ to this cluster: $N[x]=X\cup A_X$, so every pair at $x$ is satisfied, and the cost does not change. Hence $\OPT(G)\le|H|+\ell$. If $A_X=\emptyset$, $X$ is a clique component of $G$, which costs nothing, and $H$ is disjoint from it.

(N5) After the rules each class has at most $\ell+1$ vertices; the non-empty sets $A_X\subseteq V(H)$ are pairwise disjoint (Lemma~\ref{lem:R}), and $|V(H)|\le3|H|$. Classes with $A_X=\emptyset$ were deleted.
\end{proof}

\begin{remark}[Tightness of the threshold in (N4)]
For $\ell=0$ the threshold $|X|\ge2$ cannot be lowered to $|X|\ge1$: take the $4$-cycle of Proposition~\ref{prop:c4}, $H=\{(0,1,3)\}$, $X=\{2\}$. Then $(G,H,0)$ is a NO instance, but $G-2$ is the induced $P_3$ $0\,1\,3$ itself, so $(G-2,H,0)$ is a YES instance. For $\ell=1$ the threshold $|X|\ge\ell+2=3$ cannot be lowered to $|X|\ge\ell+1=2$: replace vertex $2$ of the $4$-cycle by two adjacent twins $x,x'$, i.e.\ $G$ has vertices $0,1,3,x,x'$ and edges $01,13,0x,0x',3x,3x',xx'$, and $H=\{(0,1,3)\}$ (pair-maximal, $R=X=\{x,x'\}$). Then $\OPT(G)=3$, so $(G,H,1)$ is a NO instance, whereas $\OPT(G-x)=2=|H|+1$, so $(G-x,H,1)$ is a YES instance (both by enumerating all partitions). For $\ell=2$ the same family with $|X|=3$ does not separate ($\OPT=3$ in both graphs). So the threshold is tight for $\ell\in\{0,1\}$; tightness for $\ell\ge2$ is open.
\end{remark}

\begin{remark}[The rules and the input packing]\label{rem:ext}
The input packing $H$ of CEaMP need not be pair-maximal. Extend it greedily (polynomial time) to a pair-maximal $H'\supseteq H$ with $t=|H'|-|H|$ new members, set $\ell'=\ell-t$ (answer NO if $\ell'<0$, Lemma~\ref{lem:MAX}) and apply this section to $(G,H',\ell')$. Then (N4), (N5) hold with $\ell'$ and $H'$; in particular $n\le3(\ell+2)(|H|+\ell)$ after reduction, since $|H'|\le|H|+\ell$. But $\tw(I(H'))$ may exceed $\tw(I(H))$, so these rules are a preprocessing for the parameters $(\ell,|H'|)$ or $(\ell,\tw(I(H')))$, not for $(\ell,\tw(I(H)))$ with the given $H$; and $|V(H')|$ is not bounded by $\ell$, so (N5) is not a kernel.
\end{remark}

\begin{proposition}[Partners of edits touching $R$]\label{prop:pointer}
In the setting of Corollary~\ref{cor:nf}, let $X$ be a class, $T=T_X$. Then:
\begin{enumerate}
\item[(P1)] \emph{Deletion.} If $\{x,a\}$ is an RH-deletion, $a\in A_X$, then $X$ is determined by $a$ (the unique class with $a\in A_X$): it is a unary flag at $a$ of weight $|X|$. Moreover $a$ is then cut off from every $z\in A_X\cap T$ with $z\notin Q(a)$, each such edge $az$ being a hard HH-deletion; so at least $|A_X\cap T|-q$ HH-edits accompany it.
\item[(P2)] \emph{No detachment.} If $T\cap A_X=\emptyset$ then $T=\emptyset$.
\item[(P3)] \emph{Insertion.} If $a\in T\setminus A_X$ then $T\cap A_X\ne\emptyset$, and for every $z\in T\cap A_X$ either $az$ is soft ($a\in Q(z)$, at most $q$ choices of $a$ per $z$) or $az$ is a hard non-edge, hence an inserted HH-edit.
\item[(P4)] \emph{Wide classes.} If $|A_X|\ge\ell+q$ then $T\subseteq A_X$.
\end{enumerate}
Consequently, every candidate partner $a\in V(H)$ of a class in an RH-edit is unary (P1), local to some $z$ (P3, soft case) or tied to an HH-edit (P3, second case; P1). This classifies the partners; it does not assert that the HH-edits are independent of the RH-edits.
\end{proposition}
\begin{proof}
(P1) Uniqueness is (N3). Let $a\in A_X\setminus T$ and $z\in A_X\cap T$, $z\notin Q(a)$. If $az\notin E$ then $az$ would be soft by (N3), contradicting $z\notin Q(a)$; so $az\in E$, hard, and $a\notin\Pi(X)\ni z$, so $az$ is deleted. At most $q$ vertices $z$ lie in $Q(a)$.

(P2) Suppose $T\ne\emptyset=T\cap A_X$. The cluster is $C=X\cup T$ (Theorem~\ref{thm:C1}), $a_X=|A_X\cap C|=0$, and by Lemma~\ref{lem:exch} moving $X$ out changes the cost by $|X|(|X|-|X|-|T|)=-|X||T|<0$, contradicting optimality.

(P3) $T\neq\emptyset$, so $T\cap A_X\ne\emptyset$ by (P2). Let $z\in T\cap A_X$. If $az\in E$ then $az$ is soft by (N3) ($a\in V(H)\setminus A_X$), i.e.\ $a\in Q(z)$. Otherwise either $az$ is soft (again $a\in Q(z)$) or it is a hard non-edge inside the cluster $\Pi(X)$, an insertion.

(P4) Suppose $T\not\subseteq A_X$. Put $\alpha=|A_X\cap T|$, $\delta=|A_X\setminus T|$, $\beta=|T\setminus A_X|\ge1$. By (N1) there are $|X|\beta$ RH-insertions and $|X|\delta$ RH-deletions. For $b\in T\setminus A_X$ and $z\in A_X\cap T$ with $z\notin Q(b)$, the pair $bz$ is hard, lies in a cluster, and is a non-edge by (N3); so it is a hard insertion, an HH-edit. There are at least $\beta\max(0,\alpha-q)$ of them. Hence $\ell\ge|U|\ge\beta(|X|+\max(0,\alpha-q))+|X|\delta$. If $\alpha\ge q$ this gives $\ell\ge1+(\alpha-q)+\delta$, i.e.\ $|A_X|=\alpha+\delta\le\ell+q-1$. If $\alpha<q$ it gives $\delta\le\ell-1$ and $|A_X|\le q-1+\ell-1$. Either way $|A_X|<\ell+q$.
\end{proof}

